\ssr{РЕФЕРАТ}

Расчетно-пояснительная записка содержит \pageref{LastPage}~с., \totalfigures~рис., \totaltables~табл., 18~источников, 2~приложения.

LINUX, МОДУЛЬ ЯДРА, USB-ДРАЙВЕР, СТИЛУС, ОТНОСИТЕЛЬНОЕ ПЕРЕМЕЩЕНИЕ, ПОДСИСТЕМА ВВОДА, URB.

Объект разработки --- обработка ввода стилуса графического планшета в операционной системе Linux. Предмет работы --- получение отчётов устройства через USB и их преобразование в относительные перемещения и состояния кнопок мыши.

Цель работы --- разработать загружаемый модуль ядра Linux, позволяющий использовать стилус в качестве мыши без изменения дескриптора отчёта устройства.

Сопоставлены способы обработки ввода в пространстве пользователя, на уровне HID и на уровне USB. Разработаны структура USB-драйвера, алгоритмы приёма и декодирования отчётов, масштабирование перемещений с сохранением остатков, остановка асинхронного обмена и восстановление конечной точки. Реализованы модуль на языке C и средства управления привязками интерфейсов.

Выполнена сборка для целевого ядра. По журналу событий подтверждены относительные перемещения, нажатия и отпускания обеих кнопок, а также движение при удержании левой кнопки. Численно рассмотрено влияние делителя на накопление малых перемещений.
\clearpage
